% The cloud-optimized architecture (ADR 0014) for one example query: the
% catalog request, the chunk plan built from the file URLs it returns, and the
% fan-out that sends every chunk through the load balancer to its endpoint.
%
% Chunks are colored by endpoint and numbered in plan order; the rows hidden
% by the gap hold chunks 9 to 92. Each column is as wide as its interval
% label. The dashed line is the boundary between dataretrieval-python and
% the service.
%
% Build the PNG that ADR 0014 includes, from this directory:
%   pdflatex 0014-cloud-optimized-architecture.tex
%   pdftocairo -png -r 200 -singlefile 0014-cloud-optimized-architecture.pdf \
%     0014-cloud-optimized-architecture
\documentclass[tikz,border=8pt]{standalone}
\usepackage[T1]{fontenc}
\usepackage[scaled]{helvet}
\renewcommand{\familydefault}{\sfdefault}
\usepackage{listings}
\usetikzlibrary{arrows.meta}

% one color per endpoint
\definecolor{cat}{HTML}{C2620A}   % STAC catalog
\definecolor{fil}{HTML}{2E7D32}   % Parquet files
\definecolor{api}{HTML}{1F5FA8}   % OGC API

% Python highlighting
\definecolor{pykeyword}{HTML}{CF222E}
\definecolor{pystring}{HTML}{0A3069}
\definecolor{pyname}{HTML}{8250DF}
\lstdefinestyle{python}{
  language=Python,
  basicstyle=\ttfamily\footnotesize,
  keywordstyle=\color{pykeyword},
  stringstyle=\color{pystring},
  emph={get_continuous}, emphstyle=\color{pyname},
  showstringspaces=false, columns=fullflexible, keepspaces=true,
  upquote=true, aboveskip=0pt, belowskip=0pt,
}

\begin{document}
\newsavebox{\querycode}
\begin{lrbox}{\querycode}
\begin{minipage}{4.9cm}
\begin{lstlisting}[style=python]
sites = [...]
df, md = get_continuous(
    monitoring_location_id=sites,
    time="2025-07-01/..",
)
\end{lstlisting}
\end{minipage}
\end{lrbox}

% the endpoint boxes are as wide as their widest line
\newlength{\boxwidth}
\newlength{\widthA}\settowidth{\widthA}{\sffamily\footnotesize\bfseries Parquet files}
\newlength{\widthB}\settowidth{\widthB}{\sffamily\footnotesize\bfseries STAC catalog}
\newlength{\widthC}\settowidth{\widthC}{\ttfamily\footnotesize /stac-files/}
\pgfmathsetlength{\boxwidth}{max(\widthA, \widthB, \widthC)}

\begin{tikzpicture}[
  font=\sffamily\footnotesize,
  >={Stealth[length=2.2mm]},
  chunk/.style={fill=#1!45, draw=white, line width=1pt},
  number/.style={font=\sffamily\scriptsize, text=black!75},
  axislabel/.style={font=\sffamily\scriptsize},
  column/.style={axislabel, anchor=north west, inner xsep=6pt, outer sep=0pt},
  flow/.style={->, thick},
  send/.style={->, thick, draw=#1},
  count/.style={above, font=\sffamily\scriptsize},
  stage/.style={draw=black!60, fill=black!6, rounded corners=2pt,
    minimum width=0.9cm, minimum height=6.95cm, anchor=west},
  endpoint/.style={draw=#1, very thick, rounded corners=2pt, align=left,
    fill=white, inner sep=3pt, text width=\boxwidth, anchor=west},
  region/.style={font=\sffamily\footnotesize\bfseries, text=black!65},
]

% The main flow runs along y = 2.4, through the gap in the chunk plan. The
% fan-out and the load balancer span y = -0.7 to 6.25; every stage fits
% within that height.

% ---------------------------------------------------------------- the query
\node[draw=black!50, rounded corners=2pt, inner sep=6pt, anchor=west]
  (query) at (0, 2.4) {\usebox{\querycode}};
\node[axislabel, anchor=north west, text=black!70]
  at ([yshift=-0.3cm]query.south west) {$T$: time of the last Parquet update};

% ------------------------------------------------------- the catalog request
\node[draw=cat, very thick, fill=cat!45, rounded corners=2pt,
  minimum width=0.75cm, minimum height=2.6cm, anchor=west] (catalog)
  at ([xshift=0.55cm]query.east) {};
\node[rotate=90] at (catalog.center) {catalog request};
\draw[flow] (query.east) -- (catalog.west);

% ------------------------------------------------------------ the chunk plan
% Rows: one site each, 0.5 high; sites 1 to 4, a gap, sites 47 to 50.
% Columns: the interval labels, placed first, set each column's width.
\coordinate (plan west) at ([xshift=1.55cm]catalog.east);
\draw[flow] (catalog.east) -- node[count] {file URLs} (plan west);
\coordinate (grid west) at ([xshift=1.05cm]plan west);
\node[column] (c1) at (grid west |- 0, -0.2) {2025-07-01/2026-06-26};
\node[column] (c2) at (c1.north east) {2026-06-26/$T$};
\node[column] (c3) at (c2.north east) {$T$/..};

\path (c1.west) -- (c3.east) coordinate[midway] (grid middle);
\node[anchor=south, font=\sffamily\small\bfseries] at (grid middle |- 0, 5.15)
  {Chunk plan};

\fill[chunk=api] (c3.west |- 0, 0.025) rectangle (c3.east |- 0, 4.775);
\node[number] at (c3 |- 0, 2.4) {101};
\foreach \site/\row/\top in
    {1/1/4.775, 2/2/4.275, 3/3/3.775, 4/4/3.275,
     47/47/2.025, 48/48/1.525, 49/49/1.025, 50/50/0.525} {
  \fill[chunk=fil] (c1.west |- 0, \top-0.5) rectangle (c1.east |- 0, \top);
  \fill[chunk=fil] (c2.west |- 0, \top-0.5) rectangle (c2.east |- 0, \top);
  \node[number] at (c1 |- 0, \top-0.25) {\the\numexpr 2*\row-1\relax};
  \node[number] at (c2 |- 0, \top-0.25) {\the\numexpr 2*\row\relax};
  \node[axislabel, anchor=east] at ([xshift=-0.1cm]c1.west |- 0, \top-0.25)
    {site \site};
}
\node at ([xshift=-0.5cm]c1.west |- 0, 2.4) {$\vdots$};
\node at (c1 |- 0, 2.4) {$\vdots$};
\node at (c2 |- 0, 2.4) {$\vdots$};

% x axis: one interval per column
\draw[black!50] (c1.west |- 0, -0.1) -- (c3.east |- 0, -0.1);
\foreach \edge in {c1.west, c1.east, c2.east, c3.east}
  \draw[black!50] (\edge |- 0, -0.03) -- (\edge |- 0, -0.17);

% ------------------------------------------------------------- the fan-out
% The catalog request, the file chunks, and the change chunk all run on it.
\node[stage] (fanout) at ([xshift=0.6cm]c3.east |- 0, 2.775) {};
\node[align=center, rotate=90] at (fanout.center)
  {\textbf{fan-out}\\concurrent requests};
\draw[send=cat] (catalog.north) |- (fanout.west |- 0, 5.85);
\draw[flow] (c3.east |- 0, 2.4) -- (fanout.west |- 0, 2.4);

% --------------------------------------- the boundary and the load balancer
\coordinate (boundary) at ([xshift=1.9cm]fanout.east);
\draw[black!45, dashed, thick] (boundary |- 0, -0.95) -- (boundary |- 0, 6.85);
\node[region, anchor=south east] at ([xshift=-0.12cm]boundary |- 0, 6.45)
  {dataretrieval-python};
\node[region, anchor=south west] at ([xshift=0.12cm]boundary |- 0, 6.45)
  {api.waterdata.usgs.gov};

\node[stage] (lb) at ([xshift=0.45cm]boundary |- 0, 2.775) {};
\node[rotate=90, font=\sffamily\footnotesize\bfseries] at (lb.center)
  {load balancer};

\foreach \lane/\y/\label in {cat/5.85/1 chunk, fil/2.4/100 chunks, api/0.6/1 chunk} {
  \draw[send=\lane] (fanout.east |- 0, \y) -- (lb.west |- 0, \y);
  \node[count] at ([xshift=0.95cm]fanout.east |- 0, \y) {\label};
}

% ---------------------------------------------------------------- endpoints
\node[endpoint=cat] (stac)  at ([xshift=0.5cm]lb.east |- 0, 5.85)
  {\textbf{STAC catalog}\\\texttt{/stac/}};
\node[endpoint=fil] (files) at ([xshift=0.5cm]lb.east |- 0, 2.4)
  {\textbf{Parquet files}\\\texttt{/stac-files/}};
\node[endpoint=api] (ogc)   at ([xshift=0.5cm]lb.east |- 0, 0.6)
  {\textbf{OGC API}\\\texttt{/ogcapi/}};
\foreach \lane/\y in {cat/5.85, fil/2.4, api/0.6}
  \draw[send=\lane] (lb.east |- 0, \y) -- ([xshift=0.5cm]lb.east |- 0, \y);

\end{tikzpicture}
\end{document}
